% Generated from tex/chapters/ch04/10_トピックと参考文献の対応.tex for the SWEBOK structure tree
\subsection{トピックと参考文献の対応}

この表は、SWEBOK が示すトピック対参考資料の行列を、学習用に簡略化したものである。詳しく学ぶときは、各トピックに対応する文献を読む。

\par\smallskip\noindent\refstepcounter{table}\label{tab:ch04-04}\textbf{表 \thetable: 第04章 ソフトウェア構築 表04}\par\smallskip
\begin{longtable}{L{0.46\linewidth}L{0.46\linewidth}}
\toprule
\textbf{トピック} & \textbf{主な参考資料} \\
\midrule
1. ソフトウェア構築の基礎 & McConnell, \textit{Code Complete}; Sommerville, \textit{Software Engineering}; Kim et al., \textit{The DevOps Handbook} \\
1.1 複雑さの最小化 & McConnell \\
1.2 変化を予期し受け入れる & McConnell; Sommerville; Kim et al. \\
1.3 検証しやすく構築する & McConnell \\
1.4 資産の再利用 & Sommerville \\
1.5 構築における標準の適用 & McConnell \\
2. 構築の管理 & McConnell; Sommerville; Kim et al. \\
2.1 ライフサイクルモデルにおける構築 & McConnell; Sommerville; Kim et al. \\
2.2 構築計画 & McConnell \\
2.3 構築測定 & McConnell \\
2.4 依存関係の管理 & Sommerville \\
3. 実務上の考慮 & McConnell; Sommerville; Heitkötter et al. \\
3.1 構築設計 & McConnell; Sommerville \\
3.2 構築言語 & McConnell \\
3.3 コーディング & McConnell \\
3.4 構築テスト & McConnell; Sommerville \\
3.5 構築における再利用 & Sommerville \\
3.6 構築品質 & McConnell; Sommerville \\
3.7 統合 & McConnell; Sommerville; Kim et al. \\
3.8 クロスプラットフォーム開発と移行 & Heitkötter et al. \\
4. 構築技術 & McConnell; Sommerville; Clements et al.; Gamma et al.; Mellor and Balcer; Null and Lobur; Silberschatz et al. \\
4.1 API の設計と利用 & Clements et al. \\
4.2 オブジェクト指向の実行時問題 & McConnell \\
4.3 パラメータ化、テンプレート、総称 & Gamma et al. \\
4.4 表明、契約による設計、防御的プログラミング & McConnell \\
4.5 エラー処理、例外処理、フォールトトレランス & McConnell \\
4.6 実行可能モデル & Mellor and Balcer \\
4.7 状態に基づく構築技法と表駆動構築技法 & McConnell \\
4.8 実行時設定と国際化 & McConnell \\
4.9 文法に基づく入力処理 & McConnell; Null and Lobur \\
4.10 並行処理の基本部品 & Silberschatz et al. \\
4.11 ミドルウェア & Clements et al.; Null and Lobur \\
4.12 分散およびクラウド基盤ソフトウェアの構築方法 & Sommerville; Kim et al. \\
4.13 異種システムの構築 & Null and Lobur \\
4.14 性能分析とチューニング & McConnell \\
4.15 プラットフォーム標準 & Heitkötter et al.; Null and Lobur; Silberschatz et al. \\
4.16 テストファーストプログラミング & McConnell; Sommerville \\
4.17 構築のためのフィードバックループ & Kim et al. \\
5. ソフトウェア構築ツール & McConnell; Sommerville \\
5.1 開発環境 & McConnell \\
5.2 視覚的プログラミングとローコード/ゼロコード & McConnell \\
5.3 単体テストツール & McConnell; Sommerville \\
5.4 プロファイリング、性能分析、スライシングツール & McConnell \\
\bottomrule
\end{longtable}
